\fancychapter{Conclusion}
\label{chap:Conclusion}

\textcolor{red}{last thing to write}

\section{Future Work}

The prototype developed in this thesis provides a functional platform for \ac{tsb}, but several improvements remain possible.

On the control side, additional \ac{mppt} algorithms should be implemented, tuned and compared under the same operating conditions. A dedicated irradiance measurement, used together with the existing temperature sensing, would also enable a more accurate estimation of conversion losses and could serve as an input to more advanced tracking methods.

Regarding the processing unit, a \ac{mcu} with a native 16-bit ADC would improve measurement resolution. A higher clock frequency would increase the PWM duty cycle resolution, which is particularly relevant for variable step size algorithms.

Finally, both cost and output voltage ripple could be reduced by coordinating the \acp{mppt} that operate in parallel. If the converters are interleaved, that is, operated with a controlled phase shift, the individual output ripples tend to cancel rather than add. A common output capacitor bank could then be shared among chargers, reducing the required capacitance per unit and the overall cost of the installation.

